We set out to answer a narrow question with instruments that cannot reward the easy failure modes: a semantics of justified belief that keeps alternatives open, a contract that scores behaviour rather than representation, and a benchmark that separates what happened from what was reported and from what the reports justify. In the structured setting the instruments discriminate as designed, retraction propagation is isolated from ordinary updating, and two honest negatives emerged: a learned predictor calibrates better than the hand-designed evidence model, and cue resolution---needed to keep the benchmark from being abstention-dominated---is an attack surface that provenance does not close. In the natural-language setting the primary hypothesis held by a wide margin and at lower cost, the matched-cost and raw-report variants of query-time adjudication did worse rather than better, and the sign survived a stronger adjudicator; the calibration hypothesis failed, and the corruption curve located the regime where last-write-wins overtakes justified belief. On streams assembled from real revision histories the propagation effect and the equality with a rule-based adjudicator transfer; the LLM adjudicator's much larger gap there is in good part a conflict between its priors and a declared schema, and a convention the contract failed to state, both of which we report rather than net out. On LongMemEval, a benchmark that scores recency on self-reports, a semantics that refuses to assume recency is, as it should be, penalised; what that benchmark does not contain is what this study is about. Everything needed to reproduce, contest or extend these results, including the pre-registration with its dated deviations and the cost ledger, is released.
